% ===== 原文 pp.122–131（PDF p.132–141） =====
% [原文 p.122]
\chapter{2d 拓扑场论}

\section{A-型 2d 拓扑规范理论}

本节讨论一个二维拓扑规范理论，它可以由 $d=2$、$\cN=(2,2)$ 超对称规范理论借助 $U(1)_V$ 流扭转得到 [46]。这个理论在某种程度上类似于四维的 Donaldson-Witten 理论，但要简单得多，因为它的路径积分不会获得非微扰贡献。

% [Tutor Note: p.122 术语首现。用 $U(1)_V$ 流扭转得到 A-型理论，用 $U(1)_A$ 流扭转得到 B-型理论（见 §A.2）；二者分别取自 $(2,2)$ 超对称的两个 R-对称 $U(1)$。核对方法：A-型扭转给出标量 BRST 荷并保留一个复标量 $\sigma$（本节），B-型扭转给出 $\varphi\in\Omega^1$ 且 $\cA=A+\ii\varphi$ 成为 BRST 不变的主对象（§A.2），两种理论的正定性与拓扑性均来自作用量写成 BRST-正合项这一事实。]

该理论的玻色场由一个定向二维流形 $\Sigma$ 上主 $G$-丛 $P$ 的联络 $A$，以及一个复 0-形式
\[ \sigma \in \Omega^0(\Sigma, ad\,P)\otimes\C \ . \]
费米子场为
\begin{gather*}
\beta \in \Omega^0(\Sigma, ad\,P) \\
\lambda \in \Omega^1(\Sigma, ad\,P) \\
\chi = \Omega^2(\Sigma, ad\,P) \ .
\end{gather*}
BRST 变换为
\begin{align*}
\delta_QA &= \lambda, \\
\delta_Q\sigma &= 0, \\
\delta_Q\bar\sigma &= \eta, \\
\delta_Q\beta &= 2\ii[\bar\sigma,\sigma], \\
\delta_Q\lambda &= 2\ii d_A\sigma, \\
\delta_Q\chi &= F \ .
\end{align*}

% [Tutor Note: p.122 原文两处可疑/笔误，均照原文排印。其一：费米子清单的第三行印作 $\chi=\Omega^2(\Sigma,ad\,P)$，而前两行都用 $\in$，此处应为 $\chi\in\Omega^2(\Sigma,ad\,P)$（“$\chi$ 是一个 2-形式”而非“$\chi$ 等于整个空间”）。其二：变换 $\delta_Q\bar\sigma=\eta$ 中出现了 $\eta$，但本页的费米子清单只有 $\beta,\lambda,\chi$ 三个场，$\eta$ 无定义。**判定依据**：下文说 $\delta_Q^2=2\ii\delta_\sigma$ 对除 $\chi$ 外的所有场成立，对 $\beta$ 核对给出 $\delta_Q^2\beta=2\ii([\delta_Q\bar\sigma,\sigma]+[\bar\sigma,\delta_Q\sigma])=2\ii[\delta_Q\bar\sigma,\sigma]$；要让它等于 $2\ii\delta_\sigma\beta=2\ii[\sigma,\beta]$，必须有 $[\delta_Q\bar\sigma,\sigma]=[\sigma,\beta]$，即 $\delta_Q\bar\sigma$ 必须像 $\beta$ 一样是 0-形式费米子。若按原文取 $\eta$（一个未定义的场）则该代数关系无法成立；取 $\delta_Q\bar\sigma=\beta$ 则自动成立。故此处应为 $\beta$。]

% [原文 p.123]
关系
\[ \delta_Q^2 = 2\ii\delta_\sigma \]
在除 $\chi$ 之外的所有场上离壳成立，而对 $\chi$ 有
\[ \delta_Q^2\chi = d_A\lambda \ . \]
如果存在费米型运动方程
\[ d_A\lambda = 2\ii[\chi,\sigma] \ , \]
那么关系 $\delta_Q^2=2\ii\delta_\sigma$ 在上壳成立。为了在离壳实现封闭，我们引入一个辅助玻色 2-形式 $P\in\Omega^2(\Sigma,ad\,P)$ 并重新定义 $\chi$ 的变换律：
\[ \delta_Q\chi = P, \qquad \delta_QP = 2\ii[\chi,\sigma] \ . \]

% [Tutor Note: p.123 可验证的符号核对。原文记 $\delta_\sigma$ 为“以 $\sigma$ 为参数的无穷小规范变换”。按 §2.2.1 的约定 $\delta_\epsilon A=-d_A\epsilon$，而本节又有 $\delta_Q\lambda=2\ii d_A\sigma$，于是 $\delta_Q^2A=\delta_Q\lambda=2\ii d_A\sigma=-2\ii\,\delta_\sigma A$；对 $\beta$ 的核对给出同一因子：$\delta_Q^2\beta=-2\ii\,\delta_\sigma\beta$（见上一页批注中的计算）。故若沿用 §2.2.1 的约定，关系式应为 $\delta_Q^2=-2\ii\delta_\sigma$，原文的 $2\ii\delta_\sigma$ 意味着本节取的是相反符号的 $\delta_\sigma$ 约定（或在转抄时漏了负号）。此项不必纠结：只需记住 $\delta_Q^2$ 等于“以 $\sigma$ 为参数的无穷小规范变换”，符号随约定而变；但整章内必须自洽，核对时以 $\delta_Q^2A=\delta_Q\lambda$ 为准。]

作用量的构造要保证在壳时 $P=F$。可以取一个 BRST-正合的作用量：
\[ S = \frac{1}{2e^2}\delta_Q\int_\Sigma\left(\frac{\ii}{2}\lambda\wedge\star d_A\bar\sigma + \chi\wedge\star(P-2F)\right) \ . \]

% [Tutor Note: p.123 补出省略论证。为什么“$e^2$ 只出现在 BRST-正合项的系数里”就推出它无关紧要：在取 $\delta_Q$-上同调之后，$\delta_Q(\ldots)$ 形式的项在关联函数中的贡献为零（把 $\delta_Q$ 从被积式上移走、并用可观测量 BRST 不变性），故 $e^2$ 不进入拓扑关联函数。原文“局域化到 $F=0$ 且 $\sigma$ 为常数”一句也可由此得到：令 $e^2\to0$ 时鞍点由 $F=0$、$d_A\bar\sigma=0$ 等项给出；由于 $\Sigma$ 是二维且 $F=0$，$\sigma$ 只能取常数。]

规范耦合只作为 BRST-正合项的系数出现，因此是无关紧要的。该理论的路径积分局域化在 $F=0$ 且 $\sigma$ 为常数的构型上，因此没有留下非微扰贡献的余地。

\para{对 $G=U(1)$ 的 A-型 2d 规范理论的膜范畴}

在第 6 章中我们需要理解这个 2d TFT 的膜范畴，至少是在阿贝尔情形。下面我们将论证：对 $G=U(1)$，该理论与靶空间为 $\C[2]$ 的 B-模型同构，因此其膜范畴等价于 $D^b(\C[2])$。

% [Tutor Note: p.123 记号的关键点：$\C[2]$ 不是 $\C^2$。“$[2]$”表示次数平移，$\C[2]$ 是把一维复空间放在同调次数 2 上得到的分次空间。依据就在本页下文与 p.124：Neumann 边界上 BRST 不变可观测量的代数“由 $\sigma$ 的幂张成，即 $\C[2]$ 上全纯函数的代数 $\cO$”，而 p.124 明确说玻色场的鬼数为 2；因此 $\cO=\C[\sigma]$、$\mathrm{gh}(\sigma)=2$，这正是 $\C[2]$ 上全纯函数代数的意思。若误读为 $\C^2$，“鬼数为 2”与“函数代数由单个 $\sigma$ 的多项式张成”都无法解释。另注：本页写作 $D^b(\C[2])$，而 p.124 起写作 $D^b(\operatorname{Coh}(\C[2]))$，后者才是标准写法（$\C[2]$ 上凝聚层的导出范畴），前者疑为简写或漏字。]

处理这个问题的一条进路是构造 $D^b(\operatorname{Coh}(\C[2]))$ 到膜范畴的一个嵌入。这并不能证明两个范畴等价，但确实表明前者是后者的一个满子范畴。A-型 2d 规范理论中最基本的边界条件是 Neumann 条件，它让 $A$ 与 $\sigma$ 在边界上的限制保持自由，并要求 $\star F$ 与 $\sigma$ 的法向导数在边界上为零。BRST 不变性还要求 $\star\chi$ 以及 $\star\lambda$ 在边界上的限制为零，而 $\beta$ 与 $\lambda$ 的限制则不受约束。Neumann 边界上 BRST 不变可观测量的代数由 $\sigma$ 的幂张成，也就是说，它是 $\C[2]$ 上全纯函数的代数 $\cO$。人们可以通过在边界上放置额外的自由度（它们生活在一个分次向量空间 $V$ 中）来构造更一般的边界条件。BRST 算子给出一个次数为 1 的微分 $T:V\to V$，它可以多项式地依赖于 $\sigma$。因此我们可以把膜附接到分次代数 $\cO$ 上的任意自由 DG-模 $M=(V\otimes\cO,T)$ 上。任意两个这样的膜 $M_1=(V_1\otimes\cO,T_1)$ 与 $M_2=(V_2\otimes\cO,T_2)$ 之间的态射空间是复形 $\Hom_\cO(M_1,M_2)$ 的上同调，这与范畴 $D^b(\operatorname{Coh}(\C[2]))$ 中的态射空间一致。

% [原文 p.124]
在 A-型 2d 规范理论中，也可以考虑 Dirichlet 边界条件，它在边界上令 $\sigma=0$ 并要求规范场在边界上的限制是平凡的。人们可能会猜想它对应于 $\C[2]$ 原点处的摩天楼层，而事实上可以验证：从上面考虑过的任一膜到 Dirichlet 膜的态射空间，与 $\C[2]$ 上相应的向量丛复形到该摩天楼层的态射空间一致。

\para{由与靶空间为 $\C[2]$ 的 B-模型的等价得到的膜范畴}

处理这个问题的另一条进路是构造 A-型 2d 规范理论与靶空间为 $\C[2]$ 的 B-模型之间的一个同构。从物理的观点看，两个 2d TFT $\mathbb{X}$ 与 $\mathbb{Y}$ 之间的同构就是它们之间的一个可逆拓扑缺陷线 $\mathbb{A}$。在目前的情形下，这样的缺陷线只有一个候选者。回忆一下，靶空间为 $\C[2]$ 的 B-模型有一个玻色标量 $\phi$、一个费米 1-形式 $\rho$，以及费米 0-形式 $\theta$ 与 $\xi$。BRST 变换为
\begin{align*}
\delta_Q\phi &= 0, \\
\delta_Q\bar\phi &= \eta, \\
\delta_Q\eta &= 0, \\
\delta_Q\theta &= 0, \\
\delta_Q\rho &= d\phi \ .
\end{align*}
场 $\sigma$ 的鬼数为 2，场 $\rho$ 的鬼数为 1，而场 $\eta$ 与 $\theta$ 的鬼数为 $-1$。B-模型的作用量为
\[ S = -\frac{1}{2}\delta_Q\int_\Sigma\rho\wedge\star d\bar\phi + \int_\Sigma\theta\wedge d\rho \ . \]

% [Tutor Note: p.124 原文笔误与记号说明（正文照原文译出）。（1）场清单中 “fermionic 0-forms $\theta$ and $\xi$” 的 $\xi$ 应为 $\eta$：同页 BRST 变换表只列出 $\phi,\bar\phi,\eta,\theta,\rho$，$\eta$ 已出现却缺其场清单条目，而 $\xi$ 全篇仅此一见；下文配对关系也写 “fermionic 0-forms $\beta$ and $\eta$”。故正文照原文保留 $\xi$，读者应读作 $\eta$。（2）“The field $\sigma$ has ghost number 2” 一句**不是**笔误：同页下段明确说 “鬼数为 2 的玻色子 $\sigma$ 与 $\phi$ 必须在缺陷线上等同起来”，且鬼数配对 $\sigma(2)\leftrightarrow\phi(2)$、$\beta(-1)\leftrightarrow\eta(-1)$、$\lambda(1)\leftrightarrow\rho(1)$、$\star\chi(-1)\leftrightarrow\theta(-1)$ 自洽。（3）“the fields $\rho$ has” 主谓不一致，应为 “the field $\rho$ has”。]

鬼数为 2 的玻色子 $\sigma$ 与 $\phi$ 必须在缺陷线上等同起来，允许相差一个数值因子，该因子可从作用量读出。类似地，费米 1-形式 $\lambda$ 与 $\rho$ 必须等同，费米 0-形式 $\beta$ 与 $\eta$ 也必须等同。最后，必须把 $\star\chi$ 与 $\theta$ 等同起来。于是 BRST 不变性要求 $\star F$ 在边界上为零，也就是说规范场满足 Neumann 边界条件。

注意这条缺陷线对费米子与 $\sigma$ 而言基本上就是平凡缺陷线。由于玻色 $U(1)$ 规范理论的零能扇区是平凡的，这条缺陷线的可逆性几乎是显然的。让我们更形式化地说明这一点。首先，考虑两条平行的缺陷线，它们之间夹着一条 A-型 2d 规范理论的窄条。窄条的形状是 $\R\times I$，其中 $\R$
% [原文 p.125]
参数化沿缺陷线的方向。说两条缺陷线的乘积是 B-模型中的平凡缺陷线，等价于说：在区间上、两端都取 Neumann 边界条件的 $U(1)$ 规范理论有唯一的基态。这显然是正确的，因为窄条中 $A$ 的空间分量可以用一个与时间无关的规范变换规范掉，因此区间上 $U(1)$ 规范理论的物理相空间是一个点。

% [Tutor Note: p.125 补出省略论证。上面这句“$A$ 的空间分量可被规范掉”之所以成立：在 $\R\times I$ 上取与时间无关的规范变换 $g(x)$，可令 $A_1=0$（$U(1)$ 情形 $g(x)=\exp\int^xA_1$ 总有解）；剩下的 $A_0$ 在 $1+1$ 维中只是 Lagrange 乘子、没有动力学自由度，而 Neumann 条件又排除了非平凡的 Wilson 线，故物理相空间退化为一点。这正是“两端取 Neumann 条件的区间理论有唯一基态”的含义，也说明乘积缺陷线是平凡线。]

其次，考虑相反的情形：一条 B-模型的窄条被夹在两条缺陷线之间。我们希望证明这等价于 A-型 2d 规范理论中的平凡缺陷线。窄条的形状是 $S^1\times I$。为简单起见，我们将假定挖去一条窄条后的世界面由两个连通分支组成。每个分支是一个以 $S^1$ 为边界的定向流形，而 A-型 2d 规范理论的路径积分定义了对应于 $S^1$ 的 Hilbert 空间 $V$ 中的一个向量。A-型规范理论中的任一拓扑缺陷线定义了 $V^*\otimes V\simeq\End(V)$ 中的一个元素。我们希望证明 B-模型窄条对应于 $V^*\otimes V$ 中的恒等元素。首先注意，$V$ 可以等同于规范零能自由度的 Hilbert 空间与 $\sigma$ 及费米子的零能自由度的 Hilbert 空间的张量积。如上所述，分隔 A-型 2d 规范理论与 B-模型的缺陷线在 $\sigma$ 与费米子上表现为平凡缺陷线，因此在这个扇区命题是显然的。至于规范扇区，相应的零能态 Hilbert 空间是一维的，因此 B-模型窄条正比于恒等算子。上一段的论证表明它的迹为 1，所以窄条必定是恒等算子。

% [Tutor Note: p.125 可验证的推理链。这里的论证结构是：一维 Hilbert 空间上的算子必为 $c\cdot\id$（因为 $\End(\C)=\C$），再用上一段的迹等于 1（即 $\Tr=c=1$）定出 $c=1$。这正是原文“its trace is one, so the sliver must be the identity operator”的全部内容；注意“迹”是在 $V$ 上取，而非在某个有限维子空间上取。]

\para{带规范场的 2d A-模型的膜范畴}

规范群为 $G$ 的 A-型 2d 规范理论可以耦合到一个靶空间为 $X$ 的 A-模型上，只要 $X$ 容许一个辛 $G$-作用（此处我们不写下这个带规范场的 2d A-模型）。一个自然的猜想是：相应的膜范畴是 $X$ 的 Fukaya-Floer 范畴的某种 $G$-等变版本。为使这个猜想更精确，让我们考虑一个特殊情形：$X=T^*Y$ 配上其典范辛形式。文献 [43,42] 证明了 $T^*Y$ 的 Fukaya-Floer 范畴的适当版本等价于 $Y$ 上可构造层的导出范畴。现在设 $Y$ 容许 $G$-作用。我们猜想：与靶空间为 $T^*Y$ 的 A-模型耦合的 A-型 2d 规范理论的膜范畴，等价于 $Y$ 上 $G$-等变可构造层的导出范畴。在 $Y$ 为一个点且 $G=U(1)$ 的特殊情形，后一个范畴已知等价于 $D^b(\operatorname{Coh}(\C[2]))$，这与 $G=U(1)$ 的 A-型 2d 规范理论同构于靶空间为 $\C[2]$ 的 B-模型这一事实一致。

% [Tutor Note: p.125 补出文献内容。原文只说“It was shown in [43,42]”，未说明是什么定理：那是对余切丛 $T^*Y$ 的 Fukaya-Floer 范畴与 $Y$ 上可构造层导出范畴之间的等价（A-膜 $\leftrightarrow$ 可构造层）的结果；本译本的批注仅指出这一点，不增补原文未说的内容。自洽核对：$Y=\text{点}$ 时 $T^*Y=\text{点}$、可构造层退化为分次向量空间，而 $G=U(1)$ 的等变结构给出 $\C[2]$ 上的函数代数，正好与 §A.1 的 $G=U(1)$ 结论吻合。]

% [原文 p.126]
\section{B-型 2d 拓扑规范理论}

本节讨论一个二维拓扑规范理论，它可以由 $d=2$、$\cN=(2,2)$ 超对称规范理论借助 $U(1)_A$ 流扭转得到。这个理论是 $t=\ii$ 处 GL-扭转理论的二维类比。

B-型 2d 规范理论的场是定向二维流形 $\Sigma$ 上主 $G$-丛 $P$ 的一个联络 $A$，以及
\[ \varphi \in \Omega^1(\Sigma, ad\,P) \]
和费米子场
\begin{gather*}
\beta \in \Omega^0(\Sigma, ad\,P) \\
\lambda \in \Omega^1(\Sigma, ad\,P) \\
\zeta \in \Omega^2(\Sigma, ad\,P) \ .
\end{gather*}
BRST 变换为
\begin{align*}
\delta_QA &= \lambda, \\
\delta_Q\varphi &= \ii\lambda, \\
\delta_Q\lambda &= 0, \\
\delta_Q\beta &= \ii d_A^{\star}\varphi, \\
\delta_Q\zeta &= -\ii\cF \ .
\end{align*}
这里 $\cF$ 是复化联络 $\cA=A+\ii\varphi$ 的曲率。我们将把关于 $\cA$ 的协变导数记作 $d_\cA$，把关于 $\bar\cA=A-\ii\varphi$ 的协变导数记作 $d_{\bar\cA}$，并把 $\bar\cA$ 的曲率记作 $\bar\cF$。微分算子 $d_A^{\star}$ 是 $\star d_A\star$，其中 $\star$ 是 2d Hodge 星。

% [Tutor Note: p.126 记号与可验证点。附录用五点星 $\star$ 同时表示 2d Hodge 星与 $d_A^{\star}\equiv\star d_A\star$（第 2 章 §2.5.2 用 $*$ 与 $d_A^*$，二者是同一个算子，只是符号不同，勿误认为两个不同的对象）。核对：2d 中 $\star^2$ 在 $k$-形式上的符号为 $(-1)^{k(2-k)}$，即 $\star^2=+1$（0、2-形式）与 $-1$（1-形式）；故 $d_A^{\star}$ 把 $k$-形式送到 $(k-1)$-形式——检验 $\delta_Q\beta=\ii d_A^{\star}\varphi$：左端 $\beta$ 是 0-形式，右端 $\varphi$ 是 1-形式，正合。另可验证 $\cA$ 是 BRST 不变的：$\delta_Q\cA=\delta_QA+\ii\delta_Q\varphi=\lambda+\ii(\ii\lambda)=0$，这正是“$t=\ii$”类比的来源。]

该理论有一个 $U(1)$ 鬼数对称性，在其下场 $A$、$\varphi$ 中性，场 $\lambda$ 的荷为 1，而场 $\beta$、$\zeta$ 的荷为 $-1$。BRST 变换在除 $\beta$ 之外的所有场上满足 $\delta_Q^2=0$：
\[ \delta_Q^2\beta = -d_A^{\star}\lambda \ . \]
如果使用费米型运动方程 $d_A^{\star}\lambda=0$，那么 BRST 变换在上壳时幂零。让 $\delta_Q^2=0$ 在离壳成立更为方便，因此我们引入一个辅助 0-形式 $P$ 并定义
\[ \delta_Q\beta = \ii P, \qquad \delta_QP = 0 \ . \]

% [Tutor Note: p.126 可验证的推导。$\delta_Q^2\beta=-d_A^{\star}\lambda$ 直接由 $\delta_Q\beta=\ii d_A^{\star}\varphi$ 与 $\delta_Q\varphi=\ii\lambda$ 得到：$\delta_Q(\ii d_A^{\star}\varphi)=\ii d_A^{\star}(\ii\lambda)=-d_A^{\star}\lambda$（这里用了 $\delta_Q$ 与 $\star$ 交换，因度规固定）。同理可检验鬼数：$\delta_Q$ 使鬼数 $+1$，故 $\beta$ 的荷为 $-1$、$\lambda$ 为 $1$ 与 $\delta_Q\beta=\ii P$（$P$ 中性）自洽。注意引入 $P$ 的手续与 §2.4 的 $S=\delta_QV+S_{top}$ 一样，都是把开代数补成闭代数。]
